\documentclass[10pt,a4paper]{amsart}
\usepackage{amsmath,amssymb,amsthm,graphicx}
% Retained minimal preamble for the Stacks Project chapter flat.tex (More on Flatness).
% Verbatim excerpt of the chapter preamble: Xy-pic setup (native preamble lines 14-18) and the native
% theorem declarations with the mathematical macros (native preamble lines 154-217). Copied mechanically;
% nothing added, removed or redefined. The wrapper class and the amsmath/amssymb/amsthm/graphicx base
% packages are supplied by the assembler, so the class/externaldocument part of the native preamble is omitted.
% For commutative diagrams we use Xy-pic
\usepackage[all]{xy}

% We use 2cell for 2-commutative diagrams.
\xyoption{2cell}
% Theorem environments.
%
\theoremstyle{plain}
\newtheorem{theorem}[subsection]{Theorem}
\newtheorem{proposition}[subsection]{Proposition}
\newtheorem{lemma}[subsection]{Lemma}

\theoremstyle{definition}
\newtheorem{definition}[subsection]{Definition}
\newtheorem{example}[subsection]{Example}
\newtheorem{exercise}[subsection]{Exercise}
\newtheorem{situation}[subsection]{Situation}

\theoremstyle{remark}
\newtheorem{remark}[subsection]{Remark}
\newtheorem{remarks}[subsection]{Remarks}

\numberwithin{equation}{subsection}

% Macros
%
\def\lim{\mathop{\mathrm{lim}}\nolimits}
\def\colim{\mathop{\mathrm{colim}}\nolimits}
\def\Spec{\mathop{\mathrm{Spec}}}
\def\Hom{\mathop{\mathrm{Hom}}\nolimits}
\def\Ext{\mathop{\mathrm{Ext}}\nolimits}
\def\SheafHom{\mathop{\mathcal{H}\!\mathit{om}}\nolimits}
\def\SheafExt{\mathop{\mathcal{E}\!\mathit{xt}}\nolimits}
\def\Sch{\mathit{Sch}}
\def\Mor{\mathop{\mathrm{Mor}}\nolimits}
\def\Ob{\mathop{\mathrm{Ob}}\nolimits}
\def\Sh{\mathop{\mathit{Sh}}\nolimits}
\def\NL{\mathop{N\!L}\nolimits}
\def\CH{\mathop{\mathrm{CH}}\nolimits}
\def\proetale{{pro\text{-}\acute{e}tale}}
\def\etale{{\acute{e}tale}}
\def\QCoh{\mathit{QCoh}}
\def\Ker{\mathop{\mathrm{Ker}}}
\def\Im{\mathop{\mathrm{Im}}}
\def\Coker{\mathop{\mathrm{Coker}}}
\def\Coim{\mathop{\mathrm{Coim}}}

% Boxtimes
%
\DeclareMathSymbol{\boxtimes}{\mathbin}{AMSa}{"02}

%
% Macros for moduli stacks/spaces
%
\def\QCohstack{\mathcal{QC}\!\mathit{oh}}
\def\Cohstack{\mathcal{C}\!\mathit{oh}}
\def\Spacesstack{\mathcal{S}\!\mathit{paces}}
\def\Quotfunctor{\mathrm{Quot}}
\def\Hilbfunctor{\mathrm{Hilb}}
\def\Curvesstack{\mathcal{C}\!\mathit{urves}}
\def\Polarizedstack{\mathcal{P}\!\mathit{olarized}}
\def\Complexesstack{\mathcal{C}\!\mathit{omplexes}}
% \Pic is the operator that assigns to X its picard group, usage \Pic(X)
% \Picardstack_{X/B} denotes the Picard stack of X over B
% \Picardfunctor_{X/B} denotes the Picard functor of X over B
\def\Pic{\mathop{\mathrm{Pic}}\nolimits}
\def\Picardstack{\mathcal{P}\!\mathit{ic}}
\def\Picardfunctor{\mathrm{Pic}}
\def\Deformationcategory{\mathcal{D}\!\mathit{ef}}

% --- appended 2026-10-12 (Chief): the project own definitions of the reference/slogan/history
% display environments, copied mechanically from the shared Stacks preamble lines 8-12
% (\usepackage{verbatim} plus the three \newenvironment lines, which the chapter preamble keeps
% verbatim from the project preamble). Nothing added, removed or redefined: these environments
% carry no mathematical content and the project itself typesets them as comments. Without these
% three lines a counted statement whose own text uses them fails with Environment slogan undefined.
\usepackage{verbatim}
\newenvironment{reference}{\comment}{\endcomment}
\newenvironment{slogan}{\comment}{\endcomment}
\newenvironment{history}{\comment}{\endcomment}
\usepackage{cleveref}
\crefname{originallemma}{original lemma}{original lemmas}
\Crefname{originallemma}{Original lemma}{Original lemmas}
\makeatletter
\newlabel{coherent-lemma-flat-base-change-cohomology}{{02KH}{source}{}{sourcepointer-1}{}}
\newlabel{coherent-lemma-flat-base-change-cohomology@cref}{{[originallemma][0][]02KH}{[1][0][]source}{}{}{}}
\makeatother
\crefname{originallemma}{original lemma}{original lemmas}
\Crefname{originallemma}{Original lemma}{Original lemmas}
\makeatletter
\newlabel{divisors-lemma-check-isomorphism-via-depth-and-ass}{{0AVM}{source}{}{sourcepointer-2}{}}
\newlabel{divisors-lemma-check-isomorphism-via-depth-and-ass@cref}{{[originallemma][0][]0AVM}{[1][0][]source}{}{}{}}
\makeatother
\crefname{originallemma}{original lemma}{original lemmas}
\Crefname{originallemma}{Original lemma}{Original lemmas}
\makeatletter
\newlabel{algebra-lemma-apply-grothendieck}{{0337}{source}{}{sourcepointer-3}{}}
\newlabel{algebra-lemma-apply-grothendieck@cref}{{[originallemma][0][]0337}{[1][0][]source}{}{}{}}
\makeatother
\crefname{originallemma}{original lemma}{original lemmas}
\Crefname{originallemma}{Original lemma}{Original lemmas}
\makeatletter
\newlabel{properties-lemma-criterion-normal}{{0345}{source}{}{sourcepointer-4}{}}
\newlabel{properties-lemma-criterion-normal@cref}{{[originallemma][0][]0345}{[1][0][]source}{}{}{}}
\makeatother
\crefname{originallemma}{original lemma}{original lemmas}
\Crefname{originallemma}{Original lemma}{Original lemmas}
\makeatletter
\newlabel{divisors-lemma-depth-pushforward}{{0EY0}{source}{}{sourcepointer-5}{}}
\newlabel{divisors-lemma-depth-pushforward@cref}{{[originallemma][0][]0EY0}{[1][0][]source}{}{}{}}
\makeatother
\title{A finite locally free module trivial at all codimension one points of a normal base is pulled back from the base}
\author{The Stacks Project Authors}
\date{stacks-project:flat.tex:local-snapshot-20261011; GNU FDL 1.2 or later}
\begin{document}
\maketitle
\noindent\textit{Editorial note.} All delivered mathematical text is the author's own, copied from the pinned Stacks Project chapter file flat.tex; no statement or proof step was rewritten, reordered or removed.\par\medskip
\noindent\textit{Editorial note.} The preamble is the chapter's own retained minimal preamble (sources/stacks-flat/retained-preamble.tex), a verbatim excerpt of the native Xy-pic setup and of the native theorem declarations and mathematical macros, followed by the project own definitions of the reference, slogan and history display environments (verbatim package plus three newenvironment lines) that the retained file carries at its end and that the project itself typesets as comments. No macro was invented or redefined, and the wrapper class and the amsmath/amssymb/amsthm/graphicx base packages are supplied by the assembler.\par\medskip
\noindent\textit{Editorial note.} The counted statement and proof contain no cross reference, no citation and no included file, so the delivered body is self contained.\par\medskip
\noindent\textit{Editorial note.} Canonical location of the counted proposition: Stacks tag 0EXG (sources/stacks-flat/pinned-original-source.tex line 12641).\par\medskip
\noindent\textit{Allowed standard prerequisites (editorial).}\begin{itemize}
\item Flat base change for cohomology. Cohomology of Schemes, Lemma labelled coherent-lemma-flat-base-change-cohomology, cited at proof lines 12657-12659. Transfers the isomorphism of the canonical map to the localisations at codimension one points.
\item Checking an isomorphism by depth and associated points. Divisors, Lemma labelled divisors-lemma-check-isomorphism-via-depth-and-ass, cited at proof lines 12663-12665. Reduces the isomorphism to a depth bound at points of codimension at least two.
\item Normality, depth and pushforward. Properties, Lemma labelled properties-lemma-criterion-normal, cited at proof lines 12673-12675; Divisors, Lemma labelled divisors-lemma-depth-pushforward, cited at proof lines 12675-12677; Algebra, Lemma labelled algebra-lemma-apply-grothendieck, cited at proof lines 12669-12671. Supplies depth(O\_\{S,s\}) at least 2 and the induced bound for f\_*E.
\end{itemize}\medskip
\ifdefined\hypertarget\hypertarget{sourcepointer-1}{}\fi
\noindent\textit{Source reference.} Original lemma 02KH: sources/stacks-chow/tags.txt, source lines 3372--3372. Cross-chapter standard prerequisite located in the project tag list staged in the local snapshot; invoked at proof line 12661 of the counted proof. It is retained only as a navigation pointer; every load-bearing step is inside the delivered span.\par\medskip
\ifdefined\hypertarget\hypertarget{sourcepointer-2}{}\fi
\noindent\textit{Source reference.} Original lemma 0AVM: sources/stacks-chow/tags.txt, source lines 13690--13690. Cross-chapter standard prerequisite located in the project tag list staged in the local snapshot; invoked at proof line 12665 of the counted proof. It is retained only as a navigation pointer; every load-bearing step is inside the delivered span.\par\medskip
\ifdefined\hypertarget\hypertarget{sourcepointer-3}{}\fi
\noindent\textit{Source reference.} Original lemma 0337: sources/stacks-chow/tags.txt, source lines 4046--4046. Cross-chapter standard prerequisite located in the project tag list staged in the local snapshot; invoked at proof line 12672 of the counted proof. It is retained only as a navigation pointer; every load-bearing step is inside the delivered span.\par\medskip
\ifdefined\hypertarget\hypertarget{sourcepointer-4}{}\fi
\noindent\textit{Source reference.} Original lemma 0345: sources/stacks-chow/tags.txt, source lines 4080--4080. Cross-chapter standard prerequisite located in the project tag list staged in the local snapshot; invoked at proof line 12676 of the counted proof. It is retained only as a navigation pointer; every load-bearing step is inside the delivered span.\par\medskip
\ifdefined\hypertarget\hypertarget{sourcepointer-5}{}\fi
\noindent\textit{Source reference.} Original lemma 0EY0: sources/stacks-chow/tags.txt, source lines 18675--18675. Cross-chapter standard prerequisite located in the project tag list staged in the local snapshot; invoked at proof line 12678 of the counted proof. It is retained only as a navigation pointer; every load-bearing step is inside the delivered span.\par\medskip
\begin{lemma}
\label{lemma-trivial-over-dvrs}
Let $f : X \to S$ be a morphism of schemes.
Let $\mathcal{E}$ be a finite locally free $\mathcal{O}_X$-module.
Assume
\begin{enumerate}
\item $f$ is flat and proper and $\mathcal{O}_S = f_*\mathcal{O}_X$,
\item $S$ is a normal Noetherian scheme,
\item the pullback of $\mathcal{E}$ to $X \times_S \Spec(\mathcal{O}_{S, s})$
is free for every codimension $1$ point $s \in S$.
\end{enumerate}
Then $\mathcal{E}$ is isomorphic to the pullback of a finite
locally free $\mathcal{O}_S$-module.
\end{lemma}

\begin{proof}
We will prove the canonical map
$$
\Phi : f^*f_*\mathcal{E} \longrightarrow \mathcal{E}
$$
is an isomorphism. By flat base change (Cohomology of Schemes, Lemma
\ref{coherent-lemma-flat-base-change-cohomology})
and assumptions (1) and (3) we see that
the pullback of this to $X \times_S \Spec(\mathcal{O}_{S, s})$
is an isomorphism for every codimension $1$ point $s \in S$.
By Divisors, Lemma \ref{divisors-lemma-check-isomorphism-via-depth-and-ass}
it suffices to prove that $\text{depth}((f^*f_*\mathcal{E})_x) \geq 2$
for any point $x \in X$ mapping to a point $s \in S$ of codimension $\geq 2$.
Since $f$ is flat and
$(f^*f_*\mathcal{E})_x = (f_*\mathcal{E})_s \otimes_{\mathcal{O}_{S, s}}
\mathcal{O}_{X, x}$, it suffices to prove that
$\text{depth}((f_*\mathcal{E})_s) \geq 2$, see
Algebra, Lemma \ref{algebra-lemma-apply-grothendieck}.
Since $S$ is a normal Noetherian scheme
and $\dim(\mathcal{O}_{S, s}) \geq 2$
we have $\text{depth}(\mathcal{O}_{S, s}) \geq 2$, see
Properties, Lemma \ref{properties-lemma-criterion-normal}.
Thus we get what we want from
Divisors, Lemma \ref{divisors-lemma-depth-pushforward}.
\end{proof}

\end{document}
